\section{Implementation}

\system is implemented across MAScoder templates and an SGLang-based serving stack that reuses KVFlow-style scheduling hooks.
The template side builds an agent code-base plan and serializes \prefetchhint into request metadata.
The engine side propagates the field through request parsing and scheduling, preserving observability counters such as hint count, text-bearing hint count, matched tokens, queued tokens, and exact-content match reason.
The request schema treats \prefetchhint as optional metadata.
Legacy requests therefore follow the original serving path, while MAS requests can attach an array of code-base objects without changing the user-visible prompt text.
Each hint record contains a stable code-base identifier, content signature, priority, target role, use distance, and optionally the exact text block.
The optional text field is useful for active tokenization and prefetch; the signature-only form is still useful for cache accounting and for validating that later exact-content hits correspond to template-predicted objects.

The scheduler performs two prefetch paths.
The original path prefetches KV for the next agent prefix.
The coding-aware path examines code-base hints.
Hints with exact text can be tokenized and matched against the radix/HiCache structures; hints without text are retained as scheduling metadata.
The current evaluated configuration uses device-resident cache matching and disables host-backed storage prefetch for the main table because the file-backed backend exposed a token-to-KV allocator leak in smoke tests.
This implementation choice keeps the measured path conservative.
When a hint cannot be tokenized, cannot be matched, or points to an evicted object, the scheduler records the miss and the request continues through ordinary prefill.
No user-visible output depends on a hint being honored.
The counters reported in the evaluation distinguish between hints observed by the scheduler, text-bearing hints that could be considered for active prefetch, tokens queued for code-base prefetch, and tokens actually matched by the exact-content reuse path.

The exact-content reuse layer extends cached nodes with code-anchor metadata, syntax-region type, content signature, and token span information.
The reuse decision is deliberately conservative.
Locator metadata may identify a candidate, but a mismatch in content signature rejects reuse even if the AST shape, path, or span overlaps.
When a match is allowed, the implementation computes the relative position delta between source and target spans and applies RoPE alignment before the target request decodes.
The cache node extension is local to the code-aware path.
Prefix-cache entries can still be used without code metadata, and code metadata can be absent for spans that are not part of a template object.
For spans with metadata, the runtime stores both source token offsets and normalized source text signatures.
At lookup time, the target request recomputes the normalized signature from the actual prompt content.
This avoids trusting an external manifest when the live prompt differs from the planned template.

Template metadata is represented as KV-object metadata at runtime.
Each candidate object records the producing agent, role, code-base id, content type, downstream consumers, priority, and token span.
This metadata mirrors the original AgentTemplateKV lifecycle abstraction, but \system adds code-specific fields: normalized content signature, syntax-region type, anchor signature, and match reason.
The additional fields make it possible to report whether a serving speedup came from prefix cache reuse, code-base prefetch, or exact-content segment reuse.
The benchmark harness reads these counters through the serving API response and writes per-case JSON/CSV summaries.
The paper tables are generated from those summaries rather than from hand-entered numbers.
For the 100-case serving experiment, partial runs are merged by manifest index and deduplicated before recomputing averages and percentile latency statistics.
This matters because long serving experiments can be interrupted by runtime or hardware limits; recomputing the table from unique case identifiers prevents duplicated retries from inflating hit rates.

The prototype also records failure modes explicitly.
The 32B GPTQ run is stored as a hardware-capacity limitation rather than as missing evidence, and the file-backed HiCache leak is excluded from the main performance claim.
These records are not part of the algorithm, but they keep the implementation/evaluation boundary honest: the reported gains come from device-resident exact-content reuse and scheduling metadata, not from unvalidated host load-back acceleration.
